%% USPSC-Apendice.tex
% ---
% Inicia os apêndices
% ---

\begin{apendicesenv}
% Imprime uma página indicando o início dos apêndices
\partapendices
\chapter{Comparação entre Softmax padrão e Softmax com Triton}
\label{ap:softmax_naive_triton}

Este apêndice apresenta uma comparação entre duas abordagens para cálculo de \textit{softmax} em linhas de uma matriz: uma implementação padrão em PyTorch e uma implementação otimizada com OpenAI Triton. As implementações apresentadas a seguir são adaptadas do tutorial oficial \textit{Fused Softmax} da documentação do OpenAI Triton \cite{triton2025tutorials}, e são reproduzidas aqui como material didático de apoio à discussão do Capítulo~\ref{cap_metodologia}.

\section{Implementação padrão em PyTorch}

\begin{verbatim}
import torch

import triton
import triton.language as tl
from triton.runtime import driver

DEVICE = triton.runtime.driver.active.get_active_torch_device()


def is_hip():
    return (
        triton.runtime.driver.active.get_current_target().backend == "hip"
    )


def is_cdna():
    return (
        is_hip()
        and triton.runtime.driver.active.get_current_target().arch in (
            'gfx940', 'gfx941', 'gfx942', 'gfx90a', 'gfx908'
        )
    )


def naive_softmax(x):
    """Compute row-wise softmax of X using native pytorch

    We subtract the maximum element in order to avoid overflows. Softmax is
    invariant to this shift.
    """
    # read  MN elements ; write M  elements
    x_max = x.max(dim=1)[0]
    # read MN + M elements ; write MN elements
    z = x - x_max[:, None]
    # read  MN elements ; write MN elements
    numerator = torch.exp(z)
    # read  MN elements ; write M  elements
    denominator = numerator.sum(dim=1)
    # read MN + M elements ; write MN elements
    ret = numerator / denominator[:, None]
    # in total: read 5MN + 2M elements ; wrote 3MN + 2M elements
    return ret
\end{verbatim}

\section{Implementação otimizada com Triton}

O kernel abaixo processa uma linha inteira da matriz por programa, mantendo-a
em SRAM durante todo o cálculo, de modo que cada elemento é lido da memória
global uma única vez:

\begin{verbatim}
@triton.jit
def softmax_kernel(output_ptr, input_ptr, input_row_stride,
                   output_row_stride, n_rows, n_cols,
                   BLOCK_SIZE: tl.constexpr,
                   num_stages: tl.constexpr):
    # starting row of the program
    row_start = tl.program_id(0)
    row_step = tl.num_programs(0)
    for row_idx in tl.range(row_start, n_rows, row_step,
                            num_stages=num_stages):
        # The stride represents how much we need to increase the
        # pointer to advance 1 row
        row_start_ptr = input_ptr + row_idx * input_row_stride
        # The block size is the next power of two greater than
        # n_cols, so we can fit each row in a single block
        col_offsets = tl.arange(0, BLOCK_SIZE)
        input_ptrs = row_start_ptr + col_offsets
        # Load the row into SRAM, using a mask since BLOCK_SIZE
        # may be greater than n_cols
        mask = col_offsets < n_cols
        row = tl.load(input_ptrs, mask=mask, other=-float('inf'))
        # Subtract maximum for numerical stability
        row_minus_max = row - tl.max(row, axis=0)
        numerator = tl.exp(row_minus_max)
        denominator = tl.sum(numerator, axis=0)
        softmax_output = numerator / denominator
        # Write back output to DRAM
        output_ptrs = (output_ptr + row_idx * output_row_stride
                       + col_offsets)
        tl.store(output_ptrs, softmax_output, mask=mask)
\end{verbatim}

A função auxiliar a seguir consulta as propriedades do dispositivo para
dimensionar o número de programas persistentes lançados:

\begin{verbatim}
properties = driver.active.utils.get_device_properties(DEVICE.index)
NUM_SM = properties["multiprocessor_count"]
NUM_REGS = properties["max_num_regs"]
SIZE_SMEM = properties["max_shared_mem"]
WARP_SIZE = properties["warpSize"]
target = triton.runtime.driver.active.get_current_target()
kernels = {}


def softmax(x):
    n_rows, n_cols = x.shape

    # The block size of each loop iteration is the smallest power of two
    # greater than the number of columns in `x`
    BLOCK_SIZE = triton.next_power_of_2(n_cols)

    # Another trick we can use is to ask the compiler to use more threads
    # per row by increasing the number of warps (`num_warps`)
    # over which each row is distributed.
    num_warps = 8

    # Number of software pipelining stages.
    num_stages = 4 if SIZE_SMEM > 200000 else 2

    # Allocate output
    y = torch.empty_like(x)

    # pre-compile kernel to get register usage
    # and compute thread occupancy.
    kernel = softmax_kernel.warmup(
        y, x, x.stride(0), y.stride(0), n_rows, n_cols,
        BLOCK_SIZE=BLOCK_SIZE, num_stages=num_stages,
        num_warps=num_warps, grid=(1, )
    )
    kernel._init_handles()
    n_regs = kernel.n_regs
    size_smem = kernel.metadata.shared

    if is_hip():
        # NUM_REGS represents the number of regular purpose registers.
        # On CDNA architectures this is half of all registers available.
        NUM_GPRS = NUM_REGS
        if is_cdna():
            NUM_GPRS = NUM_REGS * 2

        # MAX_NUM_THREADS represents maximum resident threads per SM.
        MAX_NUM_THREADS = properties["max_threads_per_sm"]
        max_num_waves = MAX_NUM_THREADS // WARP_SIZE
        occupancy = min(NUM_GPRS // WARP_SIZE // n_regs, max_num_waves)
        occupancy = occupancy // num_warps
    else:
        occupancy = NUM_REGS // (n_regs * WARP_SIZE * num_warps)

    occupancy = min(occupancy, SIZE_SMEM // size_smem)
    num_programs = NUM_SM * occupancy
    num_programs = min(num_programs, n_rows)

    # Create a number of persistent programs.
    kernel[(num_programs, 1, 1)](
        y, x, x.stride(0), y.stride(0), n_rows, n_cols,
        BLOCK_SIZE, num_stages
    )
    return y
\end{verbatim}

\section{Comparação técnica}

\begin{itemize}
    \item \textbf{Movimentação de memória:} a versão padrão materializa tensores intermediários para cada etapa; a versão Triton busca reduzir leituras e escritas na HBM por meio de execução em blocos e reutilização de dados em SRAM.
    \item \textbf{Paralelismo:} na implementação Triton, parâmetros como \texttt{BLOCK\_SIZE}, \texttt{num\_warps} e \texttt{num\_stages} são ajustados para aumentar ocupação e throughput.
    \item \textbf{Consciência de hardware:} o código Triton utiliza propriedades do dispositivo (SMs, registradores e memória compartilhada) para definir o número de programas persistentes, aproximando a implementação do limite arquitetural da GPU.
    \item \textbf{Objetivo prático:} enquanto a versão padrão prioriza simplicidade e legibilidade, a versão Triton prioriza eficiência computacional e redução da latência em inferência.
\end{itemize}

\chapter{Código do kernel de atenção fundido}
\label{ap:kernel_atencao}

Este apêndice apresenta o código-fonte do kernel de atenção fundido
desenvolvido neste trabalho, descrito no Capítulo~\ref{cap_metodologia}. O
kernel implementa o forward de atenção densa não causal com \textit{softmax
online}, e os acumuladores $m_i$, $l_i$ e $O_i$ são mantidos em FP32. Por
eficiência, as exponenciais são calculadas em base 2 (\texttt{tl.exp2}), com o
fator $\log_2 e$ incorporado à escala $1/\sqrt{d_k}$, transformação equivalente
a $e^x = 2^{x \log_2 e}$.

\begin{verbatim}
@triton.jit
def _attention_fwd_kernel(
    q_ptr, k_ptr, v_ptr, out_ptr,
    # strides de Q, K, V e O nos eixos [B, H, N, D]
    stride_qb, stride_qh, stride_qn, stride_qd,
    stride_kb, stride_kh, stride_kn, stride_kd,
    stride_vb, stride_vh, stride_vn, stride_vd,
    stride_ob, stride_oh, stride_on, stride_od,
    num_heads: tl.constexpr,
    seq_len: tl.constexpr,
    sm_scale_log2: tl.constexpr,
    BLOCK_M: tl.constexpr,
    BLOCK_N: tl.constexpr,
    BLOCK_D: tl.constexpr,
    HEAD_DIM: tl.constexpr,
):
    pid_m = tl.program_id(0)
    pid_bh = tl.program_id(1)

    batch_idx = pid_bh // num_heads
    head_idx = pid_bh - batch_idx * num_heads

    offs_m = pid_m * BLOCK_M + tl.arange(0, BLOCK_M)
    offs_n = tl.arange(0, BLOCK_N)
    offs_d = tl.arange(0, BLOCK_D)

    q_base = batch_idx * stride_qb + head_idx * stride_qh
    k_base = batch_idx * stride_kb + head_idx * stride_kh
    v_base = batch_idx * stride_vb + head_idx * stride_vh
    out_base = batch_idx * stride_ob + head_idx * stride_oh

    # Bloco de queries mantido em registradores durante toda a
    # varredura de keys/values
    q = tl.load(
        q_ptr + q_base + offs_m[:, None] * stride_qn
              + offs_d[None, :] * stride_qd,
        mask=(offs_m[:, None] < seq_len)
             & (offs_d[None, :] < HEAD_DIM),
        other=0.0,
    )

    # Acumuladores do softmax online (FP32)
    m_i = tl.full((BLOCK_M,), -float("inf"), tl.float32)
    l_i = tl.zeros((BLOCK_M,), tl.float32)
    acc = tl.zeros((BLOCK_M, BLOCK_D), tl.float32)

    for start_n in range(0, seq_len, BLOCK_N):
        k_offsets = start_n + offs_n
        k = tl.load(
            k_ptr + k_base + k_offsets[:, None] * stride_kn
                  + offs_d[None, :] * stride_kd,
            mask=(k_offsets[:, None] < seq_len)
                 & (offs_d[None, :] < HEAD_DIM),
            other=0.0,
        )
        v = tl.load(
            v_ptr + v_base + k_offsets[:, None] * stride_vn
                  + offs_d[None, :] * stride_vd,
            mask=(k_offsets[:, None] < seq_len)
                 & (offs_d[None, :] < HEAD_DIM),
            other=0.0,
        )

        qk = tl.dot(q, tl.trans(k)) * sm_scale_log2
        qk = tl.where(
            (offs_m[:, None] < seq_len)
            & (k_offsets[None, :] < seq_len),
            qk,
            -float("inf"),
        )

        # Atualizacao incremental do maximo, da soma e da saida
        m_ij = tl.maximum(m_i, tl.max(qk, axis=1))
        p = tl.exp2(qk - m_ij[:, None])
        alpha = tl.exp2(m_i - m_ij)
        l_i = l_i * alpha + tl.sum(p, axis=1)
        acc = acc * alpha[:, None] + tl.dot(p.to(v.dtype), v)
        m_i = m_ij

    acc = acc / l_i[:, None]

    tl.store(
        out_ptr + out_base + offs_m[:, None] * stride_on
                + offs_d[None, :] * stride_od,
        acc.to(out_ptr.dtype.element_ty),
        mask=(offs_m[:, None] < seq_len)
             & (offs_d[None, :] < HEAD_DIM),
    )
\end{verbatim}

O wrapper em Python valida as entradas, calcula a grade de lançamento com um
programa por bloco de queries e por par (lote, cabeça), e define os parâmetros
de compilação:

\begin{verbatim}
def triton_flash_attention(q, k, v, *,
                           block_m: int = 16,
                           block_n: int = 64):
    """Run the custom dense attention forward kernel."""
    _validate_qkv(q, k, v)
    if not q.is_cuda:
        raise RuntimeError(
            "triton_flash_attention requires CUDA tensors")

    q, k, v = q.contiguous(), k.contiguous(), v.contiguous()
    batch, heads, seq_len, head_dim = q.shape

    block_d = triton.next_power_of_2(head_dim)
    out = torch.empty_like(q)
    grid = (triton.cdiv(seq_len, block_m), batch * heads)
    sm_scale_log2 = (1.0 / math.sqrt(head_dim)) * _LOG2_E
    num_warps = 4 if head_dim <= 64 else 8

    _attention_fwd_kernel[grid](
        q, k, v, out,
        q.stride(0), q.stride(1), q.stride(2), q.stride(3),
        k.stride(0), k.stride(1), k.stride(2), k.stride(3),
        v.stride(0), v.stride(1), v.stride(2), v.stride(3),
        out.stride(0), out.stride(1), out.stride(2),
        out.stride(3),
        heads, seq_len, sm_scale_log2,
        BLOCK_M=block_m, BLOCK_N=block_n,
        BLOCK_D=block_d, HEAD_DIM=head_dim,
        num_warps=num_warps, num_stages=3,
    )
    return out
\end{verbatim}

\end{apendicesenv}
% ---
